\chapter{Installation et lancement}
\label{chap:installation}

\section{Prérequis}

\begin{itemize}
    \item \textbf{Système d'exploitation} : Windows 10 ou 11
        (testé), Linux et macOS (compatibilité a priori, non
        testée par les développeurs).
    \item \textbf{Python} : version 3.10 ou supérieure
        (recommandation : 3.13).
    \item \textbf{Solveur XFoil} : nécessaire uniquement pour
        l'exécution des simulations. Sous Windows, l'exécutable
        \chemin{xfoil.exe} est fourni dans le dossier
        \chemin{externaltools/xfoil/} de la distribution.
\end{itemize}

\section{Distribution exécutable Windows}

Pour les utilisateurs Windows ne souhaitant pas installer Python,
une distribution \chemin{.exe} autonome est disponible. Elle
intègre l'interpréteur Python, l'ensemble des dépendances et le
binaire XFoil.

L'application se lance en double-cliquant sur
\chemin{AifoilTools.exe}.

\begin{noteinfo}
Au premier lancement, Windows peut afficher un avertissement
\emph{SmartScreen} car l'exécutable n'est pas signé. Cliquez
sur \og Informations complémentaires \fg{} puis \og Exécuter
quand même \fg{} pour autoriser le lancement.
\end{noteinfo}

\section{Installation depuis les sources}

\subsection{Récupération du dépôt}

\begin{lstlisting}[language=bash,caption={Clonage du dépôt}]
git clone https://github.com/BenoitGFree/AifoilTools.git
cd AifoilTools
\end{lstlisting}

\subsection{Création de l'environnement Python}

\begin{lstlisting}[language=bash,caption={Création du venv et installation
des dépendances (Windows)}]
py -3 -m venv env_py3
env_py3\Scripts\activate
pip install -r requirements.txt
\end{lstlisting}

Le fichier \chemin{requirements.txt} liste les dépendances :

\begin{itemize}
    \item \code{numpy} --- calcul vectoriel ;
    \item \code{scipy} --- algorithmes numériques (interpolation,
        optimisation) ;
    \item \code{matplotlib} --- tracé graphique 2D ;
    \item \code{PySide6} --- interface graphique Qt for Python.
\end{itemize}

\subsection{Lancement de l'application}

\begin{lstlisting}[language=bash,caption={Démarrage de la GUI}]
env_py3\Scripts\python.exe run_gui.py
\end{lstlisting}

L'application s'ouvre sur la fenêtre principale présentée à la
\figref{01_main_window.png}.

\section{Vérification de l'installation}

Pour vérifier que l'environnement est correctement configuré, il
est possible d'exécuter la suite de tests unitaires :

\begin{lstlisting}[language=bash,caption={Exécution des tests}]
env_py3\Scripts\python.exe -c "import unittest; import sys; ^
sys.path.insert(0, 'sources'); sys.path.insert(0, 'tests'); ^
from test_bezier import *; ^
unittest.main(module='test_bezier', exit=False, verbosity=2)"
\end{lstlisting}

La suite complète comporte plus de 240 tests qui doivent tous
passer avec succès.

\begin{attention}
La détection automatique de l'exécutable XFoil dépend de
l'emplacement de \chemin{xfoil.exe}. L'application le cherche
dans cet ordre :
\begin{enumerate}
    \item à côté de l'exécutable de l'application (mode
        \emph{frozen} PyInstaller) ;
    \item dans le dossier \chemin{externaltools/xfoil/} ;
    \item dans le PATH système.
\end{enumerate}
Si XFoil n'est pas trouvé, l'onglet \menu{Paramétrage XFoil}
fonctionne mais le lancement d'une simulation échouera.
\end{attention}
